% Part: normal-modal-logic
% Chapter: axioms-systems
% Section: soundness

\documentclass[../../../include/open-logic-section]{subfiles}

\begin{document}

\olfileid{nml}{prf}{snd}

\olsection{可靠性}

如果所有能够被推导出的公式都是有效的，则称该推导系统是\emph{可靠的}。
当考虑模态系统时，即在推导中除了 \Ax{K} 之外还能使用某些公式 $!A_1$, \dots,~$!A_n$ 的实例，我们希望每一个可推导的公式在 $!A_1$, \dots, $!A_n$ 皆为真的任意模型中也为真。

\begin{thm}[可靠性定理]\ollabel{thm:soundness}
  如果 $!A_1$, \dots, $!A_n$ 的每个实例分别在模型类 $\mClass{C}_1$, \dots, $\mClass{C}_n$ 中有效，
  那么 $\Log{K}!A_1\dots !A_n
  \Proves !B$ 蕴含，$!B$ 在模型类 $\mClass{C}_1 \cap \dots \cap \mClass{C}_n$ 中有效。
\end{thm}

\begin{proof}
  对证明长度施归纳。为简洁起见，记 $\mClass{C} =
  \mClass{C}_1 \cap \dots \cap \mClass{C}_n$。
  \begin{enumerate}
  \item 归纳基础：若 $!B$ 有一个长度为~$1$ 的证明，则它要么是一个重言式实例，要么是~\Ax{K}\iftag{prvDiamond}{ 或~\Dual{}}{} 的实例，要么是 $!A_1$, \dots,~$!A_n$ 中某一个的实例。
      在第一种情况下，$!B$ 在 $\mClass{C}$ 中有效，因为根据 \olref[syn][tau]{prop:valid-taut}，重言式实例在\emph{任意}模型类中都有效。
      第二种情况类似，由 \olref[syn][sch]{prop:Kvalid}\iftag{prvDiamond}{ 和 \olref[syn][sch]{prop:Dual-valid}}{} 即得。
      最后在第三种情况下，由于 $!B$ 在 $\mClass{C}_i$ 中有效且 $\mClass{C} \subseteq
    \mClass{C}_i$，根据 \olref[syn][val]{prop:subset-class}，$!B$ 在 $\mClass{C}$ 中亦有效。
  \item 归纳步骤：假设 $!B$ 有一个长度为 $k>1$ 的证明。
      如果 $!B$ 是重言式实例或是 $!A_1$, \dots, $!A_n$ 中某一个的实例，可如前处理。
      现假设 $!B$ 是通过\MP{} 由前面的公式 $!C \lif !B$ 和 $!C$ 得到的。
      那么 $!C \lif !B$ 和 $!C$ 都有长度 $<k$ 的证明，根据归纳假设，它们在~$\mClass{C}$ 中有效。
      由 \olref[syn][sch]{prop:soundMP}，$!B$ 在 $\mClass{C}$ 中也有效。
      最后假设 $!B$ 是通过\Nec{} 从 $!C$ 得到的（即 $!B = \Box!C$）。
      由归纳假设，$!C$ 在 $\mClass{C}$ 中有效，再根据 \olref[syn][val]{prop:Nec-rule}，$!B$ 亦然。
  \end{enumerate}
\end{proof}

\end{document}